% Generated draft: requires mathematical and source review.
\documentclass[11pt]{article}
\usepackage[margin=25mm]{geometry}
\usepackage{fontspec}
\setmainfont{Latin Modern Roman}
\usepackage{amsmath,amssymb,mathtools,unicode-math}
\setmathfont{Latin Modern Math}
\usepackage{enumitem,needspace,xurl,hyperref}
\setlength{\parindent}{0pt}
\setlength{\parskip}{5pt}
\setlength{\emergencystretch}{3em}
\newcommand{\N}{\mathbb{N}}
\newcommand{\Z}{\mathbb{Z}}
\newcommand{\Q}{\mathbb{Q}}
\newcommand{\R}{\mathbb{R}}
\newcommand{\C}{\mathbb{C}}
\newcommand{\F}{\mathbb{F}}
\newcommand{\A}{\mathbb{A}}
\newcommand{\PP}{\mathbb{P}}
\newcommand{\E}{\mathbb{E}}
\newcommand{\eps}{\varepsilon}
\newcommand{\dd}{\,\mathrm{d}}
\begin{document}
\section*{UnsolvedMath ID: 20002440}
\textbf{Source problem code:} catalogue entry
\subsection{Definitions and mathematical statement}
The extraction begins with three questions: under what conditions do profiles for ``this equation'' exist, are they stable to multidimensional perturbations, and how should their Evans function be constructed? It suggests discrete-velocity models as a first step. The preceding equation is absent. The aid identifies this as P. Szmolyan's kinetic/Boltzmann question, rather than the subsequent Y. Li water-wave contribution. Its reconstructed target asks for collision, end-state, and speed conditions giving a positive heteroclinic profile, transverse stability criteria, and an Evans construction. The exact kinetic equation, referred to as (1.3) in the clean statement, is not supplied in this record.

For a discrete-velocity notation, write
\[
\partial_t f+A_x\partial_x f+A_y\partial_y f=Q(f),\qquad
D\bar f'=Q(\bar f),\quad D=A_x-cI,\quad z=x-ct,
\]
where $A_x,A_y$ are diagonal velocity matrices and $Q$ is a specified collision law. This notation does not fill in the missing source collision data. A collision invariant $\ell$ satisfies $\ell\cdot Q(f)=0$ and imposes the Rankine--Hugoniot condition $\ell\cdot D(f_+-f_-)=0$.

The aid's concrete four-velocity Broadwell claim uses $c=\tfrac12$ and the strictly positive profile
\[
\bar f(z)=f_-+\frac{(-2,\,2/3,\,-2,\,-2)}{16(1+e^{-z/3})},\qquad
f_-=(\tfrac12,2,1,1),\quad
f_+=(\tfrac38,\tfrac{49}{24},\tfrac78,\tfrac78).
\]
It reports all collision-invariant constraints and a one-dimensional unstable-to-stable connection on an invariant leaf. The actual four velocities and collision normalization are not displayed in the supplied summary, so those data must be recovered before verifying this as a profile for an identified model.

For perturbations $e^{\lambda t+i\xi y}v(z)$ the corresponding spectral equation, when $D$ is invertible, is
\[
D v'=\bigl(DQ(\bar f(z))-\lambda I-i\xi A_y\bigr)v.
\]
The aid reports a detailed-balance symmetrizer giving three solutions decaying at $z=-\infty$ and one decaying at $z=+\infty$ for every real $\xi$ and $\operatorname{Re}\lambda>0$. With compatible bases and a common finite $z_0$, this gives an Evans determinant
\[
E(\lambda,\xi)=\det\bigl(v^-_1(z_0),v^-_2(z_0),v^-_3(z_0),v^+_1(z_0)\bigr).
\]
Its nonvanishing is a separate stability question, not a consequence of the splitting alone. The reported rank-one far-field collision derivative $DQ(f_\pm)=a_\pm b_\pm^T$ leads, in this notation, to the dispersion determinant
\[
\det M-b_\pm^T\operatorname{adj}(M)a_\pm=0,
\qquad M=\mu D+\lambda I+i\xi A_y.
\]
The missing model data prevent specialization of these expressions here. Three collision invariants place the origin in essential spectrum at zero transverse frequency, while $\bar f'$ is a localized translation mode there. Thus the aid explicitly rejects treating the origin as a naive isolated simple Evans zero without a weighted, integrated, or renormalized construction.

The adjacent source material is attributed to Y. Li, ``Stability of travelling waves of the full water wave problem near the critical case.'' It describes travelling waves $(\eta_c,\Phi_c)=(\eta(x-ct),\Phi_c(x-ct))$ for a variety of speeds and records the multiplier $\widehat G(k)=k\tanh k$. The nonlinear displayed equations are visibly merged: the kinematic relation $\eta_t=G\Phi$ runs into the dynamic equation, whose fractions and coefficients are not reliably recoverable. This water-wave setup is preserved as adjacent context, not substituted for the missing kinetic equation or silently repaired into a full water-wave problem.
\subsection{Short English statement}
The record splices a kinetic profile/stability question with a later water-wave presentation. The aid gives a positive four-component shock and a transverse Evans framework, but the actual collision law and velocities are missing here. Essential spectrum at the translation threshold prevents a simple-zero shortcut, and no full multidimensional stability conclusion follows.
\subsection{Sources}
\begin{itemize}
\item Local UnsolvedMath registry, permanent ID $20002440$, \texttt{catalogue entry}; \texttt{aim-pdes-notes.json}, \texttt{notes[51]}, source item 4, open-problem tag.
\item American Institute of Mathematics, \emph{Stability criteria for multi-dimensional waves and patterns}, \url{https://aimath.org/WWN/multidimwaves/multidimwaves.pdf}. The aid identifies the first question with P. Szmolyan; the extracted adjacent presentation names Y. Li and the full water-wave problem near the critical case.
\item Research attempt 1: partial result, problem partially solved at run, rigor label ``proved,'' novelty confidence low. Registry classification partial progress, status partially solved, statement status reconstructed/unverified. Verification reports a splice across a page/presenter boundary but supplies no complete recovered kinetic setup. No authoritative separate citation for the model results is supplied.
\end{itemize}
\end{document}
